% ==============================================================================
% CHAPITRE 4 : TESTS ET RÉSULTATS
% ==============================================================================
\chapter{Tests et Résultats}

\begin{infobox}[Objectif du chapitre]
Ce chapitre présente la stratégie de validation du système, les métriques
de performance du modèle d'IA, les scénarios de tests fonctionnels
et une analyse critique des résultats obtenus.
\end{infobox}

% ──────────────────────────────────────────────────────────────────────────────
\section{Stratégie de Validation}
% ──────────────────────────────────────────────────────────────────────────────

La validation du système s'articule autour de trois niveaux complémentaires :

\begin{figure}[H]
\centering
\begin{tikzpicture}[
  level/.style={rectangle, rounded corners=6pt, minimum width=10cm,
                minimum height=1cm, align=center, font=\small\bfseries,
                thick, draw},
]
\node[level, fill=bleuprincipal!15, draw=bleuprincipal] at (0,3)
  {Niveau 1 : Tests Unitaires — Fonctions individuelles (nettoyage, tokenisation, modèle)};
\node[level, fill=vertprincipal!15, draw=vertprincipal] at (0,1.5)
  {Niveau 2 : Tests d'Intégration — Communication API $\leftrightarrow$ DB $\leftrightarrow$ Modèle};
\node[level, fill=orangeaccent!10, draw=orangeaccent] at (0,0)
  {Niveau 3 : Tests Fonctionnels — Scénarios utilisateur bout en bout};
\end{tikzpicture}
\caption{Pyramide de validation du système}
\label{fig:test_pyramid}
\end{figure}

% ──────────────────────────────────────────────────────────────────────────────
\section{Performances du Modèle d'IA}
% ──────────────────────────────────────────────────────────────────────────────

\subsection{Résultats d'entraînement par époque}

\begin{table}[H]
\centering
\caption{Évolution des métriques pendant l'entraînement}
\label{tab:training_history}
\renewcommand{\arraystretch}{1.5}
\begin{tabular}{ccccc}
\toprule
\rowcolor{bleuprincipal!15}
\textbf{Époque} & \textbf{Loss Train} & \textbf{Acc. Train} &
\textbf{Loss Val.} & \textbf{Acc. Val.} \\
\midrule
1 & 0.9821 & 58.6\% & 0.8743 & 64.7\% \\
\rowcolor{gray!8}
2 & 0.6234 & 74.5\% & 0.6512 & 73.5\% \\
3 & \textbf{0.3876} & \textbf{85.4\%} & \textbf{0.5981} & \textbf{76.5\%} \\
\bottomrule
\end{tabular}
\end{table}

\subsection{Métriques finales sur le jeu de test}

Le meilleur modèle (epoch 3, loss validation minimale) est évalué sur
le jeu de test final qui n'a jamais été utilisé pendant l'entraînement :

\begin{table}[H]
\centering
\caption{Rapport de classification sur le jeu de test}
\label{tab:classification_report}
\renewcommand{\arraystretch}{1.5}
\begin{tabular}{lcccc}
\toprule
\rowcolor{bleuprincipal!15}
\textbf{Classe} & \textbf{Précision} & \textbf{Rappel} &
\textbf{F1-Score} & \textbf{Support} \\
\midrule
\textcolor{rougeprincipal}{\textbf{Négatif}} (0) & 1.0000 & 0.6364 & 0.7778 & 11 \\
\rowcolor{gray!8}
\textcolor{jauneneutre!70!black}{\textbf{Neutre}} (1) & 0.6250 & 0.8333 & 0.7143 & 12 \\
\textcolor{vertprincipal}{\textbf{Positif}} (2) & 0.8462 & 1.0000 & 0.9167 & 11 \\
\midrule
\rowcolor{bleuprincipal!8}
\textbf{Accuracy globale} & \multicolumn{4}{c}{\textbf{79.41\% (27/34)}} \\
\midrule
Macro avg & 0.8237 & 0.8232 & 0.8029 & 34 \\
\rowcolor{gray!8}
Weighted avg & 0.8253 & 0.7941 & 0.7993 & 34 \\
\bottomrule
\end{tabular}
\end{table}

\subsection{Matrice de confusion}

\begin{figure}[H]
\centering
\begin{tikzpicture}
\begin{axis}[
    axis on top,
    width=8cm, height=8cm,
    xmin=-0.5, xmax=2.5,
    ymin=-0.5, ymax=2.5,
    xtick={0,1,2},
    ytick={0,1,2},
    xticklabels={Négatif, Neutre, Positif},
    yticklabels={Positif, Neutre, Négatif},
    xlabel={\textbf{Prédit}},
    ylabel={\textbf{Réel}},
    xlabel style={font=\bfseries\small},
    ylabel style={font=\bfseries\small},
    tick label style={font=\small},
    colormap={greenblue}{
      color(0)=(white) color(1)=(bleuprincipal!80)
    },
    point meta min=0, point meta max=12
]
% Matrice : [Réel=Nég: 7,3,1], [Réel=Neu: 0,10,2], [Réel=Pos: 0,0,11]
\addplot[matrix plot*, mesh/cols=3, point meta=explicit]
coordinates {
  (0,2) [7]  (1,2) [3]  (2,2) [1]
  (0,1) [0]  (1,1) [10] (2,1) [2]
  (0,0) [0]  (1,0) [0]  (2,0) [11]
};

% Annotations
\node at (axis cs:0,2) {\large\bfseries 7};
\node at (axis cs:1,2) {\large 3};
\node at (axis cs:2,2) {\large 1};
\node at (axis cs:0,1) {\large 0};
\node at (axis cs:1,1) {\large\bfseries 10};
\node at (axis cs:2,1) {\large 2};
\node at (axis cs:0,0) {\large 0};
\node at (axis cs:1,0) {\large 0};
\node at (axis cs:2,0) {\large\bfseries 11};
\end{axis}
\end{tikzpicture}
\caption{Matrice de confusion sur le jeu de test (34 exemples)}
\label{fig:confusion_matrix}
\end{figure}

\subsection{Analyse des résultats}

\textbf{Points forts observés :}
\begin{itemize}
  \item La classe \textbf{Positif} est parfaitement reconnue (rappel = 100\%,
        F1 = 0.917). Le modèle n'a jamais manqué un commentaire positif.
  \item La classe \textbf{Négatif} atteint une précision parfaite (100\%) :
        quand le modèle dit qu'un commentaire est négatif, il a toujours raison.
  \item L'accuracy globale de \textbf{79.41\%} est remarquable compte tenu
        du jeu de données limité à 225 exemples.
\end{itemize}

\textbf{Points d'amélioration :}
\begin{itemize}
  \item La confusion principale se situe entre \textbf{Négatif et Neutre}
        (3 cas) et \textbf{Neutre et Positif} (2 cas). La frontière
        entre ces classes est naturellement floue dans le langage étudiant.
  \item Le rappel de la classe Négatif (63.64\%) est le plus faible :
        certains commentaires négatifs modérés sont classés comme neutres.
\end{itemize}

% ──────────────────────────────────────────────────────────────────────────────
\section{Tests Fonctionnels de l'API}
% ──────────────────────────────────────────────────────────────────────────────

\subsection{Scénarios de test}

\begin{table}[H]
\centering
\caption{Scénarios de test fonctionnels}
\label{tab:functional_tests}
\renewcommand{\arraystretch}{1.5}
\begin{tabular}{p{0.8cm} p{5.5cm} p{3.5cm} p{2.5cm}}
\toprule
\rowcolor{bleuprincipal!15}
\textbf{ID} & \textbf{Entrée} & \textbf{Résultat attendu} & \textbf{Statut} \\
\midrule
TF01 & ``Ce cours est excellent, j'ai tout compris'' & Positif (code 200) & \textcolor{vertprincipal}{\ding{52} Passé} \\
\rowcolor{gray!8}
TF02 & ``Cours ennuyeux, je ne comprends rien'' & Négatif (code 200) & \textcolor{vertprincipal}{\ding{52} Passé} \\
TF03 & ``Le cours se passe normalement'' & Neutre (code 200) & \textcolor{vertprincipal}{\ding{52} Passé} \\
\rowcolor{gray!8}
TF04 & ``Hi'' (2 caractères) & HTTP 422 (validation) & \textcolor{vertprincipal}{\ding{52} Passé} \\
TF05 & Texte vide \texttt{``''} & HTTP 422 & \textcolor{vertprincipal}{\ding{52} Passé} \\
\rowcolor{gray!8}
TF06 & GET /api/commentaires & Liste JSON + total & \textcolor{vertprincipal}{\ding{52} Passé} \\
TF07 & GET /api/commentaires?sentiment=2 & Uniquement Positifs & \textcolor{vertprincipal}{\ding{52} Passé} \\
\rowcolor{gray!8}
TF08 & DELETE /api/commentaires/999 & HTTP 404 & \textcolor{vertprincipal}{\ding{52} Passé} \\
TF09 & GET /api/stats (base vide) & Message ``Aucun commentaire'' & \textcolor{vertprincipal}{\ding{52} Passé} \\
\rowcolor{gray!8}
TF10 & GET /health & \texttt{\{status: ok\}} & \textcolor{vertprincipal}{\ding{52} Passé} \\
\bottomrule
\end{tabular}
\end{table}

\subsection{Exemple de réponse API}

Voici un exemple de réponse JSON retournée par \texttt{POST /api/predict}
pour le commentaire ``Ce cours est excellent, le professeur explique très bien'' :

\begin{lstlisting}[style=styleJSON, caption={Exemple de réponse JSON de l'API}, label={lst:api_response}]
{
  "id": 42,
  "texte_original": "Ce cours est excellent, le professeur explique très bien",
  "matiere": "Informatique",
  "sentiment_predit": 2,
  "label_sentiment": "Positif",
  "scores": {
    "negatif": 0.0187,
    "neutre":  0.0431,
    "positif": 0.9382
  },
  "date_creation": "2026-05-12T14:23:07"
}
\end{lstlisting}

% ──────────────────────────────────────────────────────────────────────────────
\section{Tests de l'Application Mobile}
% ──────────────────────────────────────────────────────────────────────────────

\begin{table}[H]
\centering
\caption{Scénarios de test de l'application Flutter}
\label{tab:mobile_tests}
\renewcommand{\arraystretch}{1.5}
\begin{tabular}{p{0.8cm} p{5cm} p{4cm} p{2cm}}
\toprule
\rowcolor{bleuprincipal!15}
\textbf{ID} & \textbf{Action utilisateur} & \textbf{Comportement attendu} & \textbf{Statut} \\
\midrule
TM01 & Soumettre un commentaire valide & Affichage SentimentCard colorée & \textcolor{vertprincipal}{\ding{52}} \\
\rowcolor{gray!8}
TM02 & Soumettre moins de 5 caractères & Message de validation rouge & \textcolor{vertprincipal}{\ding{52}} \\
TM03 & Analyser sans serveur démarré & Message d'erreur explicite & \textcolor{vertprincipal}{\ding{52}} \\
\rowcolor{gray!8}
TM04 & Naviguer vers Historique & Chargement de la liste & \textcolor{vertprincipal}{\ding{52}} \\
TM05 & Filtrer par Positif & Seuls les commentaires positifs & \textcolor{vertprincipal}{\ding{52}} \\
\rowcolor{gray!8}
TM06 & Glisser une carte vers la gauche & Apparition du fond rouge + corbeille & \textcolor{vertprincipal}{\ding{52}} \\
TM07 & Confirmer la suppression & Commentaire supprimé, liste rafraîchie & \textcolor{vertprincipal}{\ding{52}} \\
\rowcolor{gray!8}
TM08 & Annuler la suppression & Commentaire maintenu, liste inchangée & \textcolor{vertprincipal}{\ding{52}} \\
\bottomrule
\end{tabular}
\end{table}

% ──────────────────────────────────────────────────────────────────────────────
\section{Difficultés Rencontrées et Solutions}
% ──────────────────────────────────────────────────────────────────────────────

\begin{table}[H]
\centering
\caption{Difficultés rencontrées et solutions apportées}
\label{tab:difficulties}
\renewcommand{\arraystretch}{1.5}
\begin{tabular}{p{4cm} p{5cm} p{4cm}}
\toprule
\rowcolor{bleuprincipal!15}
\textbf{Problème} & \textbf{Cause} & \textbf{Solution} \\
\midrule
\texttt{sentencepiece} manquant & CamemBERT requiert SentencePiece pour la tokenisation & \texttt{pip install sentencepiece} \\
\rowcolor{gray!8}
Peer auth PostgreSQL & Mismatch utilisateur Linux/PostgreSQL & Modifier \texttt{pg\_hba.conf} : \texttt{peer} $\rightarrow$ \texttt{md5} \\
\texttt{cmdline-tools} Android manquant & SDK mal configuré & Téléchargement manuel + configuration \texttt{JAVA\_HOME} \\
\rowcolor{gray!8}
Flutter ne joint pas l'API & Adresse localhost incorrecte sur émulateur & Utiliser \texttt{10.0.2.2} au lieu de \texttt{127.0.0.1} \\
Gradient explosion & Taux d'apprentissage trop élevé & Gradient clipping (\texttt{max\_norm=1.0}) \\
\rowcolor{gray!8}
Déséquilibre de classes initial & Dataset non équilibré & 75 exemples par classe, stratification train/test \\
\bottomrule
\end{tabular}
\end{table}

\section*{Conclusion du Chapitre}

Les tests conduits ont validé le bon fonctionnement de l'ensemble des
composants du système. Le modèle atteint une précision de 79.41\% sur le
jeu de test, avec d'excellentes performances sur les classes Positif et
Négatif. Les quelques confusions observées entre Neutre et les classes
adjacentes sont explicables par la nature ambiguë du langage étudiant et
seront améliorées avec plus de données d'entraînement.

\cleardoublepage
